\documentclass[12pt]{exam}

\usepackage[T1]{fontenc}
\usepackage[utf8]{inputenc}
\usepackage[lithuanian]{babel}
\usepackage{lmodern}
\usepackage{graphicx}
\usepackage{xcolor}
\usepackage{tikz}
\usepackage[
    a4paper,
    left=1cm,
    right=1cm,
    top=2.5cm,
    bottom=1cm
]{geometry}
\usepackage{amsmath,amssymb}
\usepackage{tasks}

% Švelnus pilkas rėmelis mokinio įrašui
\newcommand{\irasolaukas}[1]{%
    \begingroup
    \setlength{\fboxrule}{0.3pt}%
    \setlength{\fboxsep}{3pt}%
    \fcolorbox{black!35}{white}{%
        \rule{0pt}{0.45cm}\hspace{#1}%
    }%
    \endgroup
}

% Sprendimo laukas:
% argumentas nurodo linijų skaičių;
% tarpas tarp linijų – 0.7 cm.
\newcommand{\sprendimas}[1]{%
    \par\smallskip
    \noindent\textit{Sprendimas:}
    \par\nobreak
    \vspace{0.1cm}
    \noindent
    \begin{tikzpicture}
        \path[use as bounding box]
            (0,0) rectangle (\linewidth,{-0.7*#1});
        \foreach \i in {1,...,#1} {
            \draw[
                draw=black!25,
                line width=0.3pt
            ]
                (0,{-0.7*\i})
                -- (\linewidth,{-0.7*\i});
        }
    \end{tikzpicture}%
    \par\smallskip
}

% Atsakymo laukelis
\newcommand{\atsakymas}[1][9]{%
    \par\smallskip
    \noindent\textit{Atsakymas:}
    \framebox[#1cm]{\rule{0pt}{0.5cm}}
    \par
}

% Antraštė
\pagestyle{headandfoot}
\runningheadrule

\firstpageheader
{%
    \small
    \begin{tabular}{@{}l@{\hspace{0.2cm}}l@{}}
        Mokinio kodas: & \irasolaukas{2.3cm} \\[0.15cm]
        Klasė:        & \irasolaukas{2.3cm}
    \end{tabular}%
}
{Kontrolinis darbas:\\Kvadratinės lygtys}
{Dioniso variantas}

\runningheader{Kvadratinės lygtys}{}{}
\firstpagefooter{}{}{}
\runningfooter{}{}{}

% Taškų žymėjimas
\pointsdroppedatright
\bracketedpoints
\setlength{\rightpointsmargin}{1.6cm}
\renewcommand{\points}{taškai}
\renewcommand{\pointname}{Taškai:}
\marginpointname{\ taškai}

\begin{document}

\begin{questions}

\question[3]
Išspręskite lygtį $(x-3)(x+4)=2x+8$.\droppoints
\sprendimas{2}
\atsakymas[3]

\question[3]
Vienas lygties $x^2+px-21=0$ sprendinys yra $-3$. Raskite $p$ ir kitą sprendinį.\droppoints
\sprendimas{2}
\atsakymas[3]

\question[6]
Išskaidykite dauginamaisiais:\droppoints
\begin{tasks}[label-width=1.2em](3)
\task $-2x^2+3x-1$
\task $a(a+b-c)-b(c-a-b)$
\task $x^3+3x^2-16x-48$
\end{tasks}
\sprendimas{3}
\atsakymas[9]

\question[4]
Stačiakampio plotas yra $96\ \text{cm}^2$, o perimetras -- $40\ \text{cm}$. Raskite jo kraštinių ilgius.\droppoints
\sprendimas{3}
\atsakymas[3]

\question[5]
Lygties $x^2-8x+6=0$ sprendiniai yra $x_1$ ir $x_2$. Neskaičiuodami jų, sudarykite redukuotą kvadratinę lygtį, kurios sprendiniai būtų $x_1^2$ ir $x_2^2$.\droppoints
\sprendimas{5}
\atsakymas[6]

\question[5]
Išspręskite lygtį $(4+8x-x^2)(x-3)(x-5)=-120$.\droppoints
\sprendimas{6}
\atsakymas[3]

\question[7]
Raskite visas realiąsias $m$ reikšmes, su kuriomis lygties $x^2-(m+3)x+3m=0$ du skirtingus realiuosius sprendinius galima pažymėti $x_1$ ir $x_2$ taip, kad $x_1^2=x_2$.\droppoints
\sprendimas{8}
\atsakymas[4]

\question[7]
\mbox{}\droppoints
\begin{tasks}[label-width=1.2em](1)
    \task Išskaidykite dauginamaisiais $kl+2k+2l+4$.
    \task Raskite visas sveikąsias $a$ reikšmes, su kuriomis lygtis $x^2-ax+2a+2=0$ turi du skirtingus sveikuosius sprendinius.
\end{tasks}
\sprendimas{5}
\atsakymas[5]
\end{questions}

% Vertinimo lentelė puslapio apačioje
\vfill

\begin{center}
    \renewcommand{\arraystretch}{1.3}
    \begin{tabular}{|c|c|c|c|c|c|c|c|c|c|c|}
        \hline
        \textbf{Taškai}
        & 0--4 & 5--8 & 9--12 & 13--16 & 17--20
        & 21--24 & 25--29 & 30--33 & 34--37 & 38--40
        \\ \hline
        \textbf{Pažymys}
        & \textbf{1} & \textbf{2} & \textbf{3} & \textbf{4} & \textbf{5}
        & \textbf{6} & \textbf{7} & \textbf{8} & \textbf{9} & \textbf{10}
        \\ \hline
    \end{tabular}
\end{center}

\end{document}